% !TeX root = ../thesis.tex

% 中英文摘要和关键字
\begin{cnabstract} 
  以循环冷却水系统的压力缓冲需求为背景，完成冷媒膨胀罐的机械设计。设备采用立式圆筒形焊接结构，筒体内径 1800 mm、切线长度 1700 mm、有效容积 5.85 m³；主体材料 Q345R （正火）钢板，上、下标准椭圆形封头，圆筒形焊制裙座支承。设计遵循 GB/T 150—2024 及 TSG 21-2016，设计压力 $0.80$ MPa(g)（含全真空），设计温度 170 ℃。经计算，筒体与封头名义厚度确定为 14 mm；操作工况环向应力 67.7 MPa，安全系数 2.33；全真空许用外压力 0.125 MPa，水压试验应力 86.3 MPa。
 \\
 \\
 \heiti 关键词：\songti 冷媒膨胀罐；压力容器；Q355R 钢；强度计算；外压稳定性；开孔补强；水压试验
\end{cnabstract} 
  
%%======English Abstract=======%%
\begin{enabstract}
  A refrigerant expansion vessel is designed to meet the pressure-buffering requirements of a recirculating cooling water system. The vessel adopts a vertical cylindrical welded structure with an inner diameter of $1800$~mm, a tangent-to-tangent length of $1700$~mm, and an effective volume of $5.85$~m\textsuperscript{3}. It is constructed of normalized Q355R steel plates, with upper and lower standard elliptical heads and a cylindrical welded skirt support. The design complies with GB/T~150—2024 and TSG~21—2016, with a design pressure of $0.80$~MPa(g) (including full vacuum) and a design temperature of $170$~\textdegree{}C. Through calculation, the nominal thickness of the shell and heads is determined to be $14$~mm; the operating hoop stress is $67.7$~MPa with a safety factor of $2.33$; the allowable external pressure under full vacuum is $0.125$~MPa; the hydrostatic test stress is $86.3$~MPa.
  \\
  \\
  \textbf{Keywords:} Refrigerant expansion vessel; pressure vessel; Q345R steel; strength calculation; external pressure stability; opening reinforcement; hydrostatic test
\end{enabstract}
